\chapter{Submanifolds and immersions}
\label{chap:submanifolds}

In finite dimensions all reasonable notions of submanifold and immersion agree, up to some
topological phenomena such as the figure eight. In infinite dimensions they split apart, already
for Banach models and more so for locally convex ones. In this chapter we first construct
embedded submanifolds, defined by slice charts onto a closed subspace of the model, and show that
they are manifolds in their own right. Then we introduce three notions of immersion (weak,
plain, split) and the corresponding notions of submanifold given by an injective map (weakly
immersed, immersed, initial), and we prove the implications among them. Finally we show that
for a Banach domain an immersion reflects smoothness of continuous maps. The chapter closes with
a list of the infinite-dimensional subtleties; some are formalized as counterexamples in
Chapter~\ref{chap:examples}, others are only discussed. The references are
\cite[Def.~3.1.10, Def.~3.3.22, Rem.~3.3.23, Lemma~3.3.24]{GN26}, \cite[\S II.6, \S IV.4]{Nee06},
\cite{KMW24} and \cite[Cap.~2]{Lar26}.

Throughout, $E$ and $F$ are real locally convex Hausdorff spaces, not necessarily complete, and
$M$, $N$ are manifolds modeled on $E$ and $F$ (Definition~\ref{def:mb-manifold}).

\section{Embedded submanifolds}

The whole construction rests on one calculus fact: a map with values in a closed subspace is as
smooth into the subspace as it is into the ambient space.

\begin{lemma}[Corestriction to a closed subspace]
  \label{lem:sub-codrestrict}
  \lean{LocallyConvex.Analysis.ContDiffMBOn.codRestrict}
  \leanok
  Let $T \subseteq F$ be a closed subspace, $D$ a locally convex space, $V \subseteq D$ open,
  and $h \colon V \to F$ a
  $\Ck{k}$ map with $h(V) \subseteq T$. Then $h$ is $\Ck{k}$ as a map $V \to T$.
\end{lemma}

\begin{proof}
  \leanok
  Induction on $k$. The difference quotients of $h$ lie in $T$, so their limit $\dd h(x)(w)$
  lies in $T$ because $T$ is closed; the directional derivative of the corestriction is the
  corestriction of $\dd h$, and the inductive hypothesis applies to $\dd h$ on $V \times D$.
  This is \cite[Lemma~1.3.19]{GN26}.
\end{proof}

\begin{definition}[Slice chart]
  \label{def:sub-slice-chart}
  \lean{LocallyConvex.Manifolds.IsSliceChart}
  \leanok
  \uses{def:mb-manifold}
  Let $N \subseteq M$ and let $S \subseteq E$ be a subspace. A chart $\varphi \colon U \to E$ is
  a \emph{slice chart} for $N$ with model $S$ if $\varphi$ belongs to the maximal atlas of $M$
  and $\varphi(U \cap N) = \varphi(U) \cap S$, that is, $q \in N$ iff $\varphi(q) \in S$ for
  $q \in U$.
\end{definition}

\begin{definition}[Embedded submanifold]
  \label{def:embedded-submanifold}
  \lean{LocallyConvex.Manifolds.IsEmbeddedSubmanifold}
  \leanok
  \uses{def:sub-slice-chart}
  A subset $N \subseteq M$ is an \emph{embedded submanifold} modeled on the subspace
  $S \subseteq E$ if $S$ is closed and every point of $N$ lies in the domain of a slice chart
  for $N$ with model $S$.
\end{definition}

This is \cite[Def.~3.1.10]{GN26}. No complement of $S$ is required.

\begin{remark}[Slice charts and the maximal atlas]
  \label{rem:sub-maximal-atlas}
  The slice charts must be taken in the maximal atlas, i.e.\ among the charts compatible with
  the atlas of $M$, and not in the atlas itself. For instance, when $M = E$ the atlas only
  contains the identity, which is a slice chart only for subspaces.
\end{remark}

\begin{definition}[Induced charts]
  \label{def:sub-induced-chart}
  \lean{LocallyConvex.Manifolds.sliceChart, LocallyConvex.Manifolds.sliceChart_apply_coe, LocallyConvex.Manifolds.sliceChart_symm_apply_coe}
  \leanok
  \uses{def:sub-slice-chart}
  A slice chart $\varphi \colon U \to E$ induces a chart of $N$ with values in $S$, namely the
  restriction $\varphi|_{U \cap N} \colon U \cap N \to \varphi(U) \cap S$, whose inverse is the
  restriction of $\varphi^{-1}$ to $\varphi(U) \cap S$.
\end{definition}

In Lean the induced chart is a total function, and a base point of $N$ is used as a junk value
outside the target of the inverse.

\begin{lemma}[Transition maps of induced charts]
  \label{lem:sub-transition}
  \lean{LocallyConvex.Manifolds.smoothMBOn_sliceChart_transition, LocallyConvex.Manifolds.sliceChart_transition_mem}
  \leanok
  \uses{def:sub-induced-chart, def:mb-groupoid, lem:sub-codrestrict}
  If $S$ is closed, the transition map between two induced charts is smooth, and so is its
  inverse; hence it belongs to the Michal--Bastiani groupoid of $S$.
\end{lemma}

\begin{proof}
  \leanok
  \uses{lem:sub-codrestrict}
  The transition map is the restriction to $S$ of the transition map $\psi \circ \varphi^{-1}$
  of $M$, which is smooth because both charts are in the maximal atlas. It takes values in
  $S$, so it is smooth into $S$ by Lemma~\ref{lem:sub-codrestrict}. The inverse is handled by
  symmetry.
\end{proof}

\begin{theorem}[Embedded submanifolds are manifolds]
  \label{thm:sub-manifold}
  \lean{LocallyConvex.Manifolds.IsEmbeddedSubmanifold.chartedSpace, LocallyConvex.Manifolds.IsEmbeddedSubmanifold.mbManifold}
  \leanok
  \uses{def:embedded-submanifold, lem:sub-transition}
  Let $N \subseteq M$ be an embedded submanifold modeled on $S$. The induced charts form an
  atlas of $N$, with the topology induced by $M$, and with it $N$ is a manifold modeled on $S$.
\end{theorem}

\begin{proof}
  \leanok
  The atlas consists of all induced charts; at each point we choose one given by the
  definition. Compatibility is Lemma~\ref{lem:sub-transition}.
\end{proof}

\begin{proposition}[The inclusion is smooth]
  \label{prop:sub-inclusion-smooth}
  \lean{LocallyConvex.Manifolds.IsEmbeddedSubmanifold.smoothMBM_val}
  \leanok
  \uses{thm:sub-manifold, def:contmdiff}
  The inclusion $N \hookrightarrow M$ of an embedded submanifold is smooth.
\end{proposition}

\begin{proof}
  \leanok
  In an induced chart and a chart $c$ of $M$, the local representative of the inclusion is
  $(c \circ \varphi^{-1}) \circ \iota_S$, with $\iota_S \colon S \to E$ the inclusion: a smooth
  transition map of $M$ composed with a continuous linear map.
\end{proof}

The simplest example is a closed subspace of the model space itself.

\begin{proposition}[Closed subspaces]
  \label{prop:sub-subspace}
  \lean{LocallyConvex.Manifolds.isSliceChart_refl, LocallyConvex.Manifolds.isEmbeddedSubmanifold_subspace, LocallyConvex.Manifolds.isSplitSubmanifold_subspace}
  \leanok
  \uses{def:embedded-submanifold, def:split-submanifold}
  A closed subspace $S \subseteq F$ is an embedded submanifold of $F$ modeled on $S$, with the
  identity as slice chart. If $S$ is complemented, it is a split submanifold.
\end{proposition}

\begin{proof}
  \leanok
  The identity of $F$ satisfies $q \in S \iff \id(q) \in S$.
\end{proof}

\begin{remark}[Subtleties of the definition]
  \label{rem:sub-embedded-subtleties}
  When $S$ is complemented, Glöckner and Neeb call $N$ a \emph{split} submanifold
  (Definition~\ref{def:split-submanifold}). In the Banach literature (e.g.\ Bourbaki)
  ``submanifold'' usually means split, and non-split ones are called quasi-submanifolds. For
  locally convex models the splitting condition is much less useful
  \cite[Rem.~3.1.11]{GN26}. On the other hand, the usual criteria to produce submanifolds
  (regular values, split immersions) rest on the inverse or implicit function theorems, which
  are not available for general locally convex models. In this library embedded submanifolds
  are therefore produced by exhibiting slice charts by hand (see Chapter~\ref{chap:examples}).
\end{remark}

\section{The tangent space of a submanifold of a vector space}

When the ambient manifold is a locally convex space $F$ itself, the tangent space of an
embedded submanifold can be described as a subspace of $F$. This is the description used for
fields of projectors and the induced connections in the chapter on connections.

\begin{definition}[Tangent space]
  \label{def:tangent-space-submanifold}
  \lean{LocallyConvex.Manifolds.IsEmbeddedSubmanifold.tangentSpace}
  \leanok
  \uses{def:embedded-submanifold}
  Let $X \subseteq F$ be an embedded submanifold modeled on $S$ and $p \in X$. Choose a slice
  chart $\varphi$ at $p$. The tangent space of $X$ at $p$ is
  \[
    \T_p X = \{ w \in F : \dd\varphi(p)(w) \in S \}.
  \]
\end{definition}

\begin{lemma}[Independence of the chart]
  \label{lem:sub-tangent-chart}
  \lean{LocallyConvex.Manifolds.dMB_mem_iff_of_isSliceChart, LocallyConvex.Manifolds.IsEmbeddedSubmanifold.mem_tangentSpace_iff}
  \leanok
  \uses{def:tangent-space-submanifold}
  If $\varphi, \varphi'$ are slice charts at $p \in X$, then $\dd\varphi(p)(w) \in S$ iff
  $\dd\varphi'(p)(w) \in S$. Hence $\T_p X$ can be read in any slice chart, and it contains
  $\dd\varphi^{-1}(\varphi(p))(s)$ for every $s \in S$.
\end{lemma}

\begin{proof}
  \leanok
  The transition map $\tau = \varphi' \circ \varphi^{-1}$ maps the slice into the slice, so by
  closedness of $S$ its derivative at $\varphi(p)$ maps $S$ into $S$; the same holds for
  $\tau^{-1}$. The chain rule then gives the equivalence. The last claim follows since
  $\dd\varphi \circ \dd\varphi^{-1}$ is the identity.
\end{proof}

\begin{proposition}[Curves]
  \label{prop:sub-curves}
  \lean{LocallyConvex.Manifolds.IsEmbeddedSubmanifold.isCkCurve_of_smooth}
  \leanok
  \uses{thm:sub-manifold, def:ck-curve, lem:sub-codrestrict}
  A curve $\beta \colon \R \to X$ which is smooth as a curve in $F$ is a $\Ck{k}$ curve of the
  manifold $X$, for every $k$.
\end{proposition}

\begin{proof}
  \leanok
  In a slice chart $\varphi$ the local representative is the corestriction to $S$ of
  $\varphi \circ \beta$, which is smooth by Lemma~\ref{lem:sub-codrestrict}.
\end{proof}

\begin{proposition}[Tangent space of a closed subspace]
  \label{prop:sub-tangent-subspace}
  \lean{LocallyConvex.Manifolds.tangentSpace_subspace}
  \leanok
  \uses{prop:sub-subspace, def:tangent-space-submanifold}
  For a closed subspace $S \subseteq F$, $\T_p S = S$ at every $p \in S$.
\end{proposition}

\begin{proof}
  \leanok
  Read the tangent space in the identity chart, whose derivative is the identity.
\end{proof}

\section{Immersions}

Let $f \colon M \to N$ be smooth. We read its differential at $x$ in the preferred charts
$c$ at $x$ and $c'$ at $f(x)$, as the map
$v \mapsto \dd(c' \circ f \circ c^{-1})(c(x))(v)$, $E \to F$. It is continuous and linear.

\begin{lemma}[Derivatives of transition maps]
  \label{lem:sub-transition-inverse}
  \lean{LocallyConvex.Manifolds.dMB_transition_leftInverse, LocallyConvex.Manifolds.dMB_transition_injective, LocallyConvex.Manifolds.SmoothMBM.continuous_dMB_chart}
  \leanok
  \uses{def:mb-manifold, def:contmdiff}
  For charts $e, e'$ of the maximal atlas, the derivative of $e \circ e'^{-1}$ at the image
  point is a left inverse of the derivative of $e' \circ e^{-1}$; in particular the latter is
  injective. The differential of a smooth map, read in the preferred charts, is continuous.
\end{lemma}

\begin{proof}
  \leanok
  Differentiate $(e \circ e'^{-1}) \circ (e' \circ e^{-1}) = \id$ near the point, by the
  chain rule.
\end{proof}

\begin{definition}[Weak immersion]
  \label{def:weak-immersion}
  \lean{LocallyConvex.Manifolds.IsWeakImmersion}
  \leanok
  \uses{def:contmdiff}
  A smooth map $f \colon M \to N$ is a \emph{weak immersion} if its differential is injective
  at every point. No condition is imposed on its range.
\end{definition}

\begin{definition}[Immersion]
  \label{def:immersion}
  \lean{LocallyConvex.Manifolds.IsImmersion}
  \leanok
  \uses{def:contmdiff}
  A smooth map $f \colon M \to N$ is an \emph{immersion} if, at every point, its differential
  has closed range and is a topological embedding, i.e.\ a topological isomorphism onto a
  closed subspace of $F$.
\end{definition}

\begin{definition}[Split immersion]
  \label{def:split-immersion}
  \lean{LocallyConvex.Manifolds.IsSplitImmersion}
  \leanok
  \uses{def:contmdiff}
  A smooth map $f \colon M \to N$ is a \emph{split immersion} if, at every point, its
  differential has a continuous linear left inverse $r \colon F \to E$. Equivalently, the
  differential is a topological isomorphism onto a complemented subspace.
\end{definition}

\begin{proposition}[Split immersion $\Rightarrow$ immersion $\Rightarrow$ weak immersion]
  \label{prop:sub-immersion-implications}
  \lean{LocallyConvex.Manifolds.IsSplitImmersion.isImmersion, LocallyConvex.Manifolds.IsImmersion.isWeakImmersion}
  \leanok
  \uses{def:weak-immersion, def:immersion, def:split-immersion}
  Every split immersion is an immersion, and every immersion is a weak immersion.
\end{proposition}

\begin{proof}
  \leanok
  Let $D$ be the differential and $r$ a continuous left inverse. The range of $D$ is the fixed
  set of the continuous map $D \circ r$, hence closed, and $D$ is an embedding since $r$ is a
  continuous left inverse. The second implication is immediate: embeddings are injective.
\end{proof}

\begin{remark}[Dependence on charts]
  The three definitions are stated in the preferred charts. Since the derivatives of transition
  maps are linear homeomorphisms (Lemma~\ref{lem:sub-transition-inverse}), the conditions do not
  depend on this choice; the chart independence is not stated separately in Lean.
\end{remark}

\begin{proposition}[Inclusion of a closed subspace]
  \label{prop:sub-subtype-immersion}
  \lean{LocallyConvex.Manifolds.dMB_chart_subtypeL, LocallyConvex.Manifolds.isImmersion_subtypeL, LocallyConvex.Manifolds.isSplitImmersion_subtypeL_iff}
  \leanok
  \uses{def:immersion, def:split-immersion}
  The linear inclusion $S \hookrightarrow F$ of a closed subspace is an immersion. It is a
  split immersion if and only if $S$ is complemented.
\end{proposition}

\begin{proof}
  \leanok
  In the identity charts the differential is the inclusion itself, a topological embedding with
  closed range $S$. A continuous left inverse $r \colon F \to S$ gives the continuous projection
  $\iota_S \circ r$ onto $S$, and conversely.
\end{proof}

\begin{remark}[Immersions in the sense of Glöckner--Neeb]
  \label{rem:sub-gn-immersion}
  \cite[Def.~3.3.25(b)]{GN26} calls $f$ an immersion when it has, locally, the form of a
  continuous linear map with a continuous linear left inverse; this forces the range of the
  differential to be complemented. We do not use it as the definition of an immersion: it is
  much more restrictive than Definition~\ref{def:immersion}. For Banach models, an immersion with
  complemented range has this local form by the inverse function theorem; for locally convex
  models that theorem is not available.
\end{remark}

\section{Submanifolds given by an injective map}

An injective map $f \colon M \to N$ exhibits its image as a submanifold of $N$, with the
manifold structure transported from $M$. This structure is extra data: the figure eight carries
two such structures. Initial submanifolds are the exception, since their structure is unique
\cite[Lemma~3.3.24]{GN26}; this uniqueness is not formalized.

\begin{definition}[Weakly immersed and immersed submanifolds]
  \label{def:sub-immersed}
  \lean{LocallyConvex.Manifolds.IsWeaklyImmersedSubmanifold, LocallyConvex.Manifolds.IsImmersedSubmanifold}
  \leanok
  \uses{def:weak-immersion, def:immersion}
  The map $f$ exhibits its image as a \emph{weakly immersed submanifold} if it is an injective
  weak immersion, and as an \emph{immersed submanifold} if it is an injective immersion.
\end{definition}

\begin{definition}[Initial submanifold]
  \label{def:initial-submanifold}
  \lean{LocallyConvex.Manifolds.IsInitialSubmanifold}
  \leanok
  \uses{def:weak-immersion}
  The map $f$ exhibits its image as an \emph{initial submanifold} if it is an injective weak
  immersion and, for every manifold $P$ modeled on a locally convex space $H$ and every map
  $g \colon P \to M$, the map $g$ is smooth as soon as $f \circ g$ is smooth.
\end{definition}

This is \cite[Def.~3.3.22]{GN26}: smooth maps into $N$ with values in the image lift to smooth
maps into $M$. The topology of $M$ may be finer than the one induced by $N$
\cite[Rem.~3.3.23]{GN26}. In Lean the test spaces $H$ and $P$ range over the universe of the
model spaces.

\begin{proposition}
  \label{prop:sub-weakly-immersed}
  \lean{LocallyConvex.Manifolds.IsImmersedSubmanifold.isWeaklyImmersedSubmanifold, LocallyConvex.Manifolds.IsInitialSubmanifold.isWeaklyImmersedSubmanifold}
  \leanok
  \uses{def:sub-immersed, def:initial-submanifold, prop:sub-immersion-implications}
  Immersed submanifolds and initial submanifolds are weakly immersed.
\end{proposition}

\begin{proof}
  \leanok
  The first claim is Proposition~\ref{prop:sub-immersion-implications}; the second holds by
  definition.
\end{proof}

\section{Embedded and split submanifolds as immersions}

\begin{definition}[Split submanifold]
  \label{def:split-submanifold}
  \lean{LocallyConvex.Manifolds.IsSplitSubmanifold}
  \leanok
  \uses{def:embedded-submanifold}
  An embedded submanifold $X \subseteq M$ modeled on $S$ is \emph{split} if $S$ is
  complemented in $E$ \cite[Def.~3.1.10]{GN26}.
\end{definition}

\begin{lemma}[Differential of the inclusion]
  \label{lem:sub-dmb-inclusion}
  \lean{LocallyConvex.Manifolds.IsEmbeddedSubmanifold.exists_dMB_chart_val}
  \leanok
  \uses{thm:sub-manifold, lem:sub-transition-inverse}
  Let $X \subseteq M$ be an embedded submanifold modeled on $S$ and $p \in X$. In the preferred
  charts, the differential of the inclusion at $p$ is the restriction to $S$ of a continuous
  linear automorphism $L$ of $E$.
\end{lemma}

\begin{proof}
  \leanok
  \uses{lem:sub-transition-inverse}
  Near the point, the local representative is $(c \circ \varphi^{-1}) \circ \iota_S$, where
  $\varphi$ is the slice chart and $c$ the preferred chart of $M$. Take $L$ the derivative of
  $c \circ \varphi^{-1}$; its inverse is the derivative of $\varphi \circ c^{-1}$ by
  Lemma~\ref{lem:sub-transition-inverse}.
\end{proof}

\begin{theorem}[The inclusion is an immersion]
  \label{thm:sub-inclusion-immersion}
  \lean{LocallyConvex.Manifolds.IsEmbeddedSubmanifold.isImmersion_val, LocallyConvex.Manifolds.IsEmbeddedSubmanifold.isImmersedSubmanifold_val}
  \leanok
  \uses{lem:sub-dmb-inclusion, def:immersion, def:sub-immersed, prop:sub-inclusion-smooth}
  The inclusion of an embedded submanifold is an immersion; hence it exhibits $X$ as an
  immersed submanifold.
\end{theorem}

\begin{proof}
  \leanok
  With $L$ as in Lemma~\ref{lem:sub-dmb-inclusion}, the range $L(S) = (L^{-1})^{-1}(S)$ is
  closed, and $L|_S$ is an embedding since $L^{-1}$ is a continuous left inverse.
\end{proof}

\begin{theorem}[Split submanifolds]
  \label{thm:sub-inclusion-split}
  \lean{LocallyConvex.Manifolds.IsSplitSubmanifold.isSplitImmersion_val}
  \leanok
  \uses{lem:sub-dmb-inclusion, def:split-submanifold, def:split-immersion}
  The inclusion of a split submanifold is a split immersion.
\end{theorem}

\begin{proof}
  \leanok
  A continuous projection of $E$ onto $S$, composed with $L^{-1}$, is a continuous left inverse
  of $L|_S$.
\end{proof}

\begin{theorem}[The inclusion is initial]
  \label{thm:sub-inclusion-initial}
  \lean{LocallyConvex.Manifolds.IsEmbeddedSubmanifold.isInitialSubmanifold_val}
  \leanok
  \uses{thm:sub-inclusion-immersion, def:initial-submanifold, lem:sub-codrestrict}
  The inclusion of an embedded submanifold exhibits it as an initial submanifold: a map into
  $X$ is smooth as soon as it is smooth into $M$.
\end{theorem}

\begin{proof}
  \leanok
  \uses{lem:sub-codrestrict}
  In a slice chart $\varphi$, the local representative of a map $g \colon P \to X$ is the
  corestriction to $S$ of the local representative of $g$ as a map into $M$, read in
  $\varphi$. The latter is smooth, and so is the former by Lemma~\ref{lem:sub-codrestrict}.
\end{proof}

Summing up, for a subset $X \subseteq M$ and an injective map $f \colon M \to N$,
\[
  \text{split} \Rightarrow \text{embedded} \Rightarrow
  \text{immersed and initial} \Rightarrow \text{weakly immersed}.
\]
None of these implications can be reversed, and initial and immersed submanifolds are not
comparable; see Section~\ref{sec:sub-subtleties}.

\section{Immersions with a Banach domain}

In this section $E$ and $F$ are normed spaces and $E$ is complete, so that $M$ is a Banach
manifold; $F$ need not be complete. Recall that a manifold modeled on an incomplete normed space
is not a Banach manifold.

\begin{theorem}[Immersions reflect smoothness]
  \label{thm:sub-banach-reflect}
  \lean{LocallyConvex.Manifolds.IsImmersion.smoothMBM_of_comp}
  \leanok
  \uses{def:immersion, def:contmdiff}
  Let $E$ be a Banach space, $F$ a normed space, and $f \colon M \to N$ an immersion. Let $P$
  be a manifold modeled on a locally convex space $H$ and $g \colon P \to M$ continuous. If
  $f \circ g$ is smooth, then $g$ is smooth.
\end{theorem}

\begin{proof}
  \leanok
  Work in charts, with $J$ the local representative of $f$ at $g(x)$. For Banach spaces,
  Michal--Bastiani smooth maps are Fréchet smooth, so $J'$ is continuous in norm. The
  differential $J'(y_0)$ is an embedding, hence bounded below, and a bound with half the
  constant persists on a neighborhood $V$ of $y_0$. On the open set where the local
  representative $\gamma$ of $g$ lands in $V$, $J \circ \gamma$ is $\Ck{k}$, and the calculus
  result \texttt{LocallyConvex.Analysis.ContDiffMBOn.of\_comp\_of\_boundedBelow} of \cite{LCA} gives that
  $\gamma$ is $\Ck{k}$: the derivative of $\gamma$ is obtained as the limit of difference
  quotients, which are Cauchy by the lower bound.
\end{proof}

\begin{corollary}[Topologically embedded immersions are initial]
  \label{cor:sub-banach-initial}
  \lean{LocallyConvex.Manifolds.IsImmersedSubmanifold.isInitialSubmanifold}
  \leanok
  \uses{thm:sub-banach-reflect, def:initial-submanifold, def:sub-immersed,
    prop:sub-immersion-implications}
  Let $E$ be a Banach space and $F$ a normed space. An injective immersion $f \colon M \to N$
  which is a homeomorphism onto its image exhibits its image as an initial submanifold.
\end{corollary}

\begin{proof}
  \leanok
  If $f \circ g$ is smooth, it is continuous, so $g$ is continuous because $f$ is a topological
  embedding. Then $g$ is smooth by Theorem~\ref{thm:sub-banach-reflect}.
\end{proof}

\begin{remark}[Scope of the Banach results]
  \label{rem:sub-banach-scope}
  No inverse function theorem is used, and no complement of the range of the differential is
  needed. On the other hand, the image of a non-split immersion need not have slice charts:
  whether such an immersion, when it is a homeomorphism onto its image, is an embedded
  submanifold is open. Completeness of $E$ is used, to produce the derivative of the lift as a
  limit of Cauchy difference quotients; completeness of $F$ is not needed. For Fréchet or more
  general locally convex model spaces the statement is open, since the lower bound of the
  differential need not persist near a point. Continuity of $g$ cannot be dropped: the figure
  eight (Remark~\ref{rem:sub-initial-embedded}) is an injective immersion that is not initial,
  because a smooth curve in the plane can lift discontinuously.
\end{remark}

\section{Subtleties in infinite dimensions}
\label{sec:sub-subtleties}

We collect the phenomena that separate the notions above. For each we say whether it is
formalized (in Chapter~\ref{chap:examples}) or only discussed.

\begin{remark}[Weak immersions versus immersions; formalized]
  \label{rem:sub-weak-vs-immersion}
  The differential can be injective without having closed range, already between Banach
  spaces: the inclusion $\ell^1 \to \ell^2$ is continuous, linear and injective, and its range
  is dense and not closed. Hence it is a weak immersion and not an immersion. This is formalized
  (\texttt{not\_isImmersion\_lpInclusion}, Chapter~\ref{chap:examples}).
\end{remark}

\begin{remark}[Immersions versus split immersions; formalized]
  \label{rem:sub-immersion-vs-split}
  A closed subspace need not be complemented, already in Banach spaces: by Phillips' lemma
  \cite{Phi40}, $c_0$ has no complement in $\ell^\infty$. So $c_0 \subseteq \ell^\infty$ is an
  embedded submanifold that is not split, and the inclusion is an immersion that is not a split
  immersion (compare Proposition~\ref{prop:sub-subtype-immersion}). This is formalized, with
  Whitley's proof \cite{Whi66} (\texttt{Phillips.not\_isSplitImmersion\_c0},
  Chapter~\ref{chap:examples}).
\end{remark}

\begin{remark}[The universal property alone; not formalized]
  \label{rem:sub-joris}
  The universal property of Definition~\ref{def:initial-submanifold} alone does not give an
  injective differential. By Joris' theorem \cite{Jor82} ($g^2$ and $g^3$ smooth imply $g$
  smooth), the cusp $t \mapsto (t^2, t^3)$ has it, and its differential vanishes at $0$. This is
  why injectivity of the differential is part of the definition \cite[p.~336]{GN26};
  Karshon, Miyamoto and Watts call the universal property alone a \emph{diffeological
  submanifold} \cite{KMW24}.
\end{remark}

\begin{remark}[Initial versus embedded; the figure eight is formalized]
  \label{rem:sub-initial-embedded}
  Initial does not mean embedded: the dense line in the torus is initial (not formalized).
  Embedded does not follow from injective immersion either. The figure eight
  $t \mapsto (\sin t, \sin 2t)$, $t \in (-\pi, \pi)$, is an injectively immersed curve that is
  not initial: the curve $s \mapsto (\sin s, -\sin 2s)$ is smooth in $\R^2$, stays in the figure
  eight, and crosses the origin along the other branch, so read in the parameter it jumps from
  $\pi$ to $0$ to $-\pi$. This is formalized, with the parameter $\theta = 2\arctan t$
  (\texttt{Lemniscate.isImmersedSubmanifold\_lemniscate} and
  \texttt{Lemniscate.not\_isInitialSubmanifold\_lemniscate}, Chapter~\ref{chap:examples}).
\end{remark}

\begin{remark}[Mechanisms that destroy initiality]
  \label{rem:sub-mechanisms}
  For an injective weak immersion, smooth maps into the ambient manifold that land in the image
  may fail to lift:
  \begin{enumerate}
    \item continuously, because two pieces of the image accumulate on each other (the figure
      eight; this occurs in finite dimensions and is formalized);
    \item continuously, because there are uncountably many pieces ($\R \times \R_{\mathrm{discrete}}
      \to \R^2$); this is excluded for second countable domains, which is why Glöckner and Neeb
      assume separability in \cite[Thm.~9.2.15]{GN26} (not formalized);
    \item continuously, because the map is not even locally a topological embedding (the linear
      inclusion $\ell^1 \to \ell^2$; infinite dimensions only; that it is not an immersion is
      formalized, the failure of initiality is only discussed);
    \item continuously, but not smoothly, even when the map is a homeomorphism onto its image
      (infinite dimensions only; not formalized). Let $A$ act diagonally on the $n$-th block of
      $\ell^2(\N, \R^2)$ by a scalar $a_n \le e^{-n}$ and let
      $j \colon \ell^2 \to \ell^2 \times \R$, $j(x) = (Ax, \|x\|^2)$. Its differential
      $v \mapsto (Av, 2\langle x, v\rangle)$ is injective, with non-closed range at $0$, and $j$
      is a homeomorphism onto its image (weak convergence plus convergence of norms). The curve
      $y(t) = \bigoplus_n n^{-2}(\cos nt, \sin nt)$ has constant norm and $A \circ y$ is smooth,
      so $j \circ y$ is smooth; but $y$ is $\Ck{1}$ and not $\Ck{2}$ in $\ell^2$. So $j$ is
      topologically embedded, weakly immersed and not initial. In finite dimensions a weak
      immersion whose lifts are continuous is initial, so this phenomenon has no
      finite-dimensional analogue. Compare Corollary~\ref{cor:sub-banach-initial}, where the
      closed range of the differential rules it out.
  \end{enumerate}
\end{remark}

\begin{remark}[Unitary orbits; not formalized]
  \label{rem:sub-unitary-orbits}
  Let $A = i\sum_n \lambda_n p_n$ be a skew-adjoint operator on a Hilbert space $\mathcal{H}$,
  with distinct real $\lambda_n$ and pairwise orthogonal projections $p_n$ (of any rank) summing
  to $1$. Its stabilizer in $U(\mathcal{H})$ is the block-diagonal group
  $K = \bigoplus_n U(p_n\mathcal{H})$, a split Lie subgroup via the conditional expectation
  $X \mapsto \sum_n p_n X p_n$. The orbit map $U(\mathcal{H})/K \to B(\mathcal{H})$,
  $uK \mapsto uAu^*$, is a weak immersion with differential $X \mapsto [X, A]$, and it is
  initial if and only if there are finitely many $p_n$.
  If there are finitely many $p_n$, the spectrum is finite and the orbit is a closed embedded
  split submanifold \cite{BL23}; with two projections of infinite rank it is a Grassmannian.
  If there are infinitely many $p_n$, the spectrum is infinite, so $[\,\cdot\,, A]$ does not
  have closed range \cite{BL23}, and the orbit is not initial. Indeed, for orbits of Banach--Lie
  group actions with split stabilizer, initiality forces closed range: along
  $s \mapsto \exp(s\xi)\cdot A$ every derivative is bounded by a constant times
  $\|[\xi, A]\|$, and bump curves along directions with $\|[\xi, A]\|$ tiny are smooth in
  $B(\mathcal{H})$ while their lifts to $U(\mathcal{H})/K$ are continuous but not $\Ck{1}$.
\end{remark}

\begin{remark}[Initial does not imply immersion; not formalized]
  \label{rem:sub-initial-not-immersion}
  Let $A$ be a compact injective diagonal operator on $\ell^2$ and
  \[
    j \colon \ell^2 \to \ell^2 \times (\ell^2 \otimes \ell^2) \times
      (\ell^2 \otimes \ell^2 \otimes \ell^2), \qquad
    j(x) = (Ax, x \otimes x, x \otimes x \otimes x),
  \]
  with Hilbert tensor products. Then $j$ is injective, its differential
  $v \mapsto (Av, x \otimes v + v \otimes x, \dots)$ is injective, and at $0$ its range
  $A(\ell^2) \times 0 \times 0$ is not closed. Nevertheless $j$ is a homeomorphism onto its
  image, the image is closed, and $j$ is initial with respect to all test manifolds modeled on
  Fréchet spaces. The argument, which is only discussed, runs as follows. Let $f \colon P \to
  \ell^2$ with $j \circ f$ smooth.
  \begin{enumerate}
    \item Along a smooth curve $c$, each coordinate $\langle f \circ c, e\rangle$ has smooth
      square and cube, namely $\langle (f \circ c) \otimes (f \circ c), e \otimes e\rangle$ and
      its cubic analogue; so it is smooth by Joris' theorem \cite{Jor82}. A curve in a Banach
      space all of whose scalar components are smooth is smooth, by uniform boundedness and
      Taylor's formula.
    \item Smoothness along curves gives Michal--Bastiani smoothness when $P$ is modeled on a
      Fréchet space. First, the directional derivatives $\dd^k f(p)(v, \dots, v)$ are locally
      bounded: otherwise there are $p_m \to p_0$ and small $w_m$ with
      $\|\dd^k f(p_m)(w_m, \dots, w_m)\| > m$. With $t_m = 2^{-m}$, $r_m = 2^{-m-2}$ and a bump
      $\beta$ supported in $[-1,1]$, equal to $1$ on $[-\tfrac12, \tfrac12]$, the curve
      $c(t) = p_0 + \sum_m \beta((t - t_m)/r_m)\,(p_m - p_0 + (t - t_m)w_m)$ is smooth if
      $p_m - p_0$ and $w_m$ decay fast enough in every seminorm. Near $t_m$ it is affine,
      $c(t) = p_m + (t - t_m)w_m$, so $(f \circ c)^{(k)}(t_m) = \dd^k f(p_m)(w_m, \dots, w_m)$
      is unbounded, a contradiction. Local bounds then give multilinearity (Boman's theorem on
      finite-dimensional subspaces) and joint continuity of the derivatives (mean value
      inequality).
  \end{enumerate}
  So for this $j$ the closed-range condition fails at $0$ although $j$ is initial, and initial
  and immersed submanifolds are not comparable. Whether $j$ is initial with respect to test
  manifolds modeled on non-metrizable spaces is open: the argument above uses sequences.
\end{remark}

\begin{remark}[Initial Lie subgroups; not formalized]
  \label{rem:sub-initial-lie}
  Glöckner and Neeb call an initial Lie subgroup \emph{strictly initial} when the induced map
  of Lie algebras is injective \cite[Def.~9.1.5]{GN26}, a condition that Neeb includes in the
  definition \cite{Nee06}. Not every subgroup carries an initial Lie subgroup structure:
  $\ell^\infty \subseteq \R^{\N}$ does not \cite[Ex.~9.1.8]{GN26}.
\end{remark}
